\chapter*{How to Read This Book}
\addcontentsline{toc}{chapter}{How to Read This Book}
\markboth{HOW TO READ THIS BOOK}{HOW TO READ THIS BOOK}

\begin{flushright}
\itshape Lasciate ogne fretta, voi ch'intrate.\\[2pt]
\upshape\small not quite Dante, \emph{Inferno}, canto III
\end{flushright}

Over the gate of Hell, Dante wrote: abandon all hope, ye who enter here. This
book asks for less: abandon all haste.

This is a dense book. It is packed with names, dates, formulas, numbers and
code, and almost every page leans on something said before it. It assumes you
will stop, go back a few pages, check a number or run an example, and come
back. It was written for curious people: the ones who want to know why a thing
is the way it is, who would rather spend an evening on one page than skim
twenty, and who are not satisfied until the idea and the history behind it both
make sense.

One project runs through the whole book: a program that simulates a chemical
plant. It moves the plant forward in time, lets other programs read and change
it, and answers questions about it. The plant is the Tennessee Eastman process,
which J.\,J.~Downs and E.\,F.~Vogel, of the Eastman Chemical Company, published
in 1993 as a challenge problem for process control.

Building that program touches five fields that usually live apart, with their
own textbooks, their own notation and their own people. The book has one part
for each, and each part answers one question the project forces on you.

The first part is about differential equations and dynamical systems, and the
question is what the plant is. A plant is a set of quantities that change at
rates that depend on each other, which is a system of differential equations.
Before computing anything, you want to know how that system behaves: where it
settles, what pushes it away, and whether it comes back.

The second part is about numerical methods and algorithms: how a computer moves
those equations forward. Almost no real system has a formula for its solution,
so the computer takes small steps, and every step carries an error. The method
decides whether that error stays small or swallows the answer.

The third part is about software engineering: how to build the program that does
all this and keep it standing as it grows. Design, patterns and code conventions
came from projects that collapsed under their own weight, and a simulator
collapses the same way.

The fourth part is about optimization: once the model runs, where should the
plant operate? Every operating point has a cost, and some are forbidden.
Optimization finds the cheapest one that is allowed.

The fifth part is about control theory: how to keep the plant there.
Disturbances push the plant away from the point the optimizer chose, and a
controller has to push it back without making things worse.

The parts are in the order in which the questions depend on each other, but the
chapters are written in the order the project reaches each subject. A part can
grow while the one before it is still thin, and when a chapter needs something
from another part, it says where to find it. None of the parts replaces a
textbook in its field. Each one goes as far as the project needs and stops
there.

\section*{Every idea has a parent}

No method in this book showed up because it was elegant. Poincaré stopped
looking for formulas and started looking at the shape of the solutions because
most differential equations, the three-body problem among them, have no formula
to find. The word \emph{stiff} entered numerical analysis in 1952, when Curtiss
and Hirschfelder found that their chemical kinetics equations forced the usual
methods into absurdly small steps. A NATO conference put the name
\emph{software engineering} on the table in 1968 because large programs kept
arriving late, over budget and full of errors. Dantzig wrote the simplex method
in 1947 because the US Air Force was planning its training, supply and
deployment by hand. Minorsky wrote down the PID controller in 1922 after
watching how helmsmen steer a ship. Each idea fixed something its predecessor
got wrong and brought costs of its own. If you know the lineage, you know when
the idea applies. If you only know the rule, you will apply it where it does not
belong.

\section*{Five kinds of boxes}

Stories, concepts, questions, reflections and words are kept apart, so you
always know which one you are reading.

\begin{history}{year \textperiodcentered{} place}
An amber box tells a story: who did it, when, and what forced the idea into
existence. The people, the dates and the problems in it are real. No invented
incidents.
\end{history}

\begin{concept}{name}
A blue box defines an idea you will need again. It always ends the same way:

\textbf{What it fixes:} the problem that made the idea worth having.

\textbf{What it costs:} what you give up by using it.
\end{concept}

\begin{question}{the question you were embarrassed to ask}
A teal box asks the question that occurs to anyone reading the paragraph above
it and that nobody asks out loud, and then answers it.
\end{question}

\begin{reflection}{the thing that should bother you}
A purple box stops on a statement that sounds absurd, or that the text just said
in passing as if it were obvious, and works out what it means and why it is
true.
\end{reflection}

\begin{etymology}
A gray box says where a word comes from. The origin of a word is not its
definition, but it often shows what the people who coined it were looking at.
\end{etymology}

The text between the boxes is the argument that connects them.

\section*{Code}

The examples are in Rust, and sometimes in Python. Source code sits on a light
gray background and what a command printed sits on a dark one. Every example was
run before it went into the book, and the output shown is the output it gave.

\section*{Diagrams}

Colors in diagrams are roles, not decoration. A color means the same thing in
every figure of the book.

\begin{center}
\begin{tikzpicture}
  \matrix[column sep=6mm, row sep=4mm, nodes={anchor=center}] {
    \node[hf/source, minimum width=26mm] {source}; &
    \node[hf/note, anchor=west, text width=95mm, align=left] {what is being simulated: a plant, a model, its equations}; \\
    \node[hf/reader, minimum width=26mm] {reader}; &
    \node[hf/note, anchor=west, text width=95mm, align=left] {a program that uses the model: a solver, an optimizer, a controller, a screen}; \\
    \node[hf/store, minimum width=26mm] {store}; &
    \node[hf/note, anchor=west, text width=95mm, align=left] {where results are kept: a time series, a snapshot, a file}; \\
    \node[hf/bus, minimum width=26mm] {middle}; &
    \node[hf/note, anchor=west, text width=95mm, align=left] {something that sits between: an adapter, a protocol, a standard}; \\
    \draw[hf/ask] (-13mm,0) -- (13mm,0); &
    \node[hf/note, anchor=west, text width=95mm, align=left] {asks: a request, a call}; \\
    \draw[hf/answer] (-13mm,0) -- (13mm,0); &
    \node[hf/note, anchor=west, text width=95mm, align=left] {answers a request}; \\
    \draw[hf/push] (-13mm,0) -- (13mm,0); &
    \node[hf/note, anchor=west, text width=95mm, align=left] {pushes: sent without being asked}; \\
    \draw[hf/save] (-13mm,0) -- (13mm,0); &
    \node[hf/note, anchor=west, text width=95mm, align=left] {keeps: written to storage}; \\
    \draw[hf/pain, dashed] (-13mm,0) -- (13mm,0); &
    \node[hf/note, anchor=west, text width=95mm, align=left, text=hfAlert] {a cost, a pain, or something forbidden}; \\
  };
\end{tikzpicture}
\end{center}

\section*{Exact facts, firm opinions}

This is not a scientific book, and it does not pretend to be one.

The history is checked. When the book says who did something and when, that is
what happened.

The opinions are another matter. When a method survives only because the
textbooks keep teaching it, the book says so. When a design pattern adds
nothing, the book says it adds nothing, not that it is ``of limited value in
some contexts''. Sometimes that will be unfair or imprecise, and we accept that.

A reader who leaves with a position can argue with it, refine it, or throw it
away. A reader who leaves with ``it depends'' has nothing to hold. Between being
precise and leaving you with a position, this book prefers to give you a
position.
